\section{Исходная программа}

\subsection{Родительский процесс}
\begin{center}
\lstinputlisting[
    language=C++,
    caption={Исходный код родительского процесса},
    captionpos=b
]{inc/parent.cpp}
\end{center}

\subsection{Дочерний процесс}
\begin{center}
\lstinputlisting[
    language=C++,
    caption={Исходный код дочернего процесса},
    captionpos=b
]{inc/child.cpp}
\end{center}

\subsection{Общие функции}
\begin{center}
\lstinputlisting[
    language=C++,
    caption={Функции разбора числа и проверки простоты},
    captionpos=b
]{inc/common.cpp}
\end{center}

\subsection{Заголовочные файлы}

\lstinputlisting[
    language=C++,
    caption={Общие объявления},
    captionpos=b
]{inc/include/common.hpp}

\lstinputlisting[
    language=C++,
    caption={Коды сообщений протокола обмена},
    captionpos=b
]{inc/include/protocol.hpp}

\subsection{Сборка}

\lstinputlisting[
    caption={Файл конфигурации CMake},
    captionpos=b
]{CMakeLists.txt}

\subsection{Проверка системных вызовов}

Для просмотра системных вызовов Linux-версии применяется утилита
\texttt{strace}. Команда запуска:

\begin{lstlisting}[language={}]
printf 'result.txt\n4\n9\n7\n' | strace -f \
-e trace=process,pipe,dup2,read,write,close \
./build/bin/parent
\end{lstlisting}

В выводе должны присутствовать системные вызовы \texttt{pipe()},
\texttt{fork()}, \texttt{dup2()}, \texttt{execve()}, \texttt{read()},
\texttt{write()}, \texttt{close()} и \texttt{wait4()}. Вызов
\texttt{execve()} соответствует запуску дочерней программы, а
\texttt{wait4()} является системной реализацией библиотечной функции
\texttt{waitpid()}.